\documentclass[12pt]{article}
\usepackage[utf8]{inputenc}
\usepackage[T1]{fontenc}
\usepackage{hyperref}
\title{Karkar-Yuri Language Resources for K-12 Curriculum Development}
\author{Compiled by Claude Sonnet 5.5 (claude-sonnet-5-5) for Andr\'as Kornai}
\date{}
\begin{document}
\maketitle

\section{Introduction}
Karkar-Yuri (ISO 639-3 \texttt{yuj}, Glottocode kark1258) is a Pauwasi language spoken by roughly a thousand people in Green River Rural LLG, Sandaun Province, Papua New Guinea, near the Indonesian border. It is the only Eastern Pauwasi language in Papua New Guinea. Published materials appear scarce. The known SIL works by Dorothy Price and Veda Rigden seem to be unpublished.

\section{Primary Educational Resources}
I found no primers, textbooks or school programmes for Karkar-Yuri.

\section{Linguistic and Reference Materials}
\begin{itemize}
\item Wikipedia, ``Karkar language'': overview and classification. \url{https://en.wikipedia.org/wiki/Karkar_language}
\item Wikipedia, ``Pauwasi languages''. \url{https://en.wikipedia.org/wiki/Pauwasi_languages}
\item Wikipedia, ``East Pauwasi languages''. \url{https://en.wikipedia.org/wiki/East_Pauwasi_languages}
\item New Guinea World: East Pauwasi River. Its bibliography lists work by Dorothy Price and Veda Rigden, including a 1987 study of discourse and plurals and a 1988 dictionary. Both are marked as unobtained there. \url{https://newguineaworld.linguistik.uzh.ch/families/pauwasi-river/east-pauwasi-river/start}
\item ASJP entry for Karkar-Yuri, which cites Rigden's unpublished ``Karkar Grammar Essentials'' held at Ukarumpa. \url{https://asjp.clld.org/languages/KARKAR_YURI}
\item Endangered Languages Project bibliography for Karkar-Yuri. \url{https://endangeredlanguages.com/lang/10680/bibliography}
\end{itemize}

\section{Religious and Community Resources}
I found none in my searches.

\section{Digital and Audio Resources}
I found no audio, video or digitized texts.

\section{Curriculum Development Context}
There is no known school programme and no accessible published grammar or dictionary. The SIL manuscripts (Price and Rigden) would be the first thing to ask about, via SIL PNG. A curriculum would depend on new community-led documentation.

\end{document}